\diarytitle{0120}{累次積分の積分順序の交換（三角形領域、e^{x^2} の積分）}

連続関数 $f(x,y)$ が有界閉領域 $D$ 上で定義されているとき、$D$ を
\[
D = \{(x,y) \mid a \le x \le b,\ \varphi_1(x) \le y \le \varphi_2(x)\}
\]
のように $x$ を先に固定して $y$ の範囲を表す「縦線集合」として書ける場合、重積分は
\[
\iint_{D} f(x,y)\,dA = \int_{a}^{b}\left(\int_{\varphi_1(x)}^{\varphi_2(x)} f(x,y)\,dy\right)dx
\]
と累次積分に書き換えられる（Fubini の定理）。同じ領域 $D$ を「横線集合」
\[
D = \{(x,y) \mid c \le y \le d,\ \psi_1(y) \le x \le \psi_2(y)\}
\]
として表せば、積分順序を交換して
\[
\iint_{D} f(x,y)\,dA = \int_{c}^{d}\left(\int_{\psi_1(y)}^{\psi_2(y)} f(x,y)\,dx\right)dy
\]
としてもよい。両者は同じ重積分の値を与える。$x$ 積分が初等関数で書けない場合（例えば $e^{x^2}$）でも、順序を交換することで積分可能になることがある。

積分領域は $0 \le y \le 1,\ y \le x \le 1$ である。これは
\[
0 \le x \le 1,\qquad 0 \le y \le x
\]
と同じ領域を表す（$x$-$y$ 平面上で直線 $y=x$ の下側、$0\le x\le1$ の三角形領域）。したがって積分順序を交換すると
\[
\int_{0}^{1}\int_{y}^{1} e^{x^{2}}\,dx\,dy = \int_{0}^{1}\int_{0}^{x} e^{x^{2}}\,dy\,dx
\]
内側の積分は $x$ を定数とみて
\[
\int_{0}^{x} e^{x^{2}}\,dy = x\,e^{x^{2}}
\]
よって
\[
\int_{0}^{1} x\,e^{x^{2}}\,dx = \left[\frac{1}{2}e^{x^{2}}\right]_{0}^{1} = \frac{e-1}{2}
\]
となるから
\[
\boxed{\; \int_{0}^{1}\int_{y}^{1} e^{x^{2}}\,dx\,dy = \frac{e-1}{2} \;}
\]
（検算: $\dfrac{d}{dx}\left(\dfrac{1}{2}e^{x^2}\right) = x e^{x^2}$ なので原始関数の計算は正しい。また元の積分順序のままでは $\int e^{x^2}dx$ が初等関数で表せないため、順序交換が本質的である。）
